% concatenated from Supplementary_R1.tex
\documentclass[12pt]{article}
\usepackage{amsmath,amssymb,amsthm}
\usepackage{bm}
\usepackage{graphicx}
\usepackage{natbib}
\usepackage{geometry}
\usepackage{xr}
\geometry{margin=1in}

\theoremstyle{remark}
\theoremstyle{plain}
\newtheorem{axiom}{Axiom}
\newtheorem{claim}[axiom]{Claim}
\newtheorem{theorem}{Theorem}[section]
\newtheorem{lemma}[theorem]{Lemma}
\newtheorem{proposition}{Proposition}[section]
\newtheorem{remark}{Remark}[section]
%\theoremstyle{remark}{Remark}[section]
\newtheorem{definition}[theorem]{Definition}
\newtheorem{example}{Example}[section]
\newtheorem*{fact}{Fact}
\newtheorem{rem}{Remark}
\newtheorem{corollary}{Proposition}
\newtheorem{assumption}{Assumption}


\DeclareMathOperator{\tr}{tr}
\DeclareMathOperator{\Var}{Var}
\DeclareMathOperator{\Cov}{Cov}
\DeclareMathOperator{\vecop}{vec}



% Link to main article
\externaldocument{TestsApproach_M}  % Make sure TestsApproach.aux is available when compiling

\title{Supplementary Material for \\
	“A Relativity-Based Framework for Statistical Testing Guided by the Independence of Ancillary Statistics: Methodology and Nonparametric Illustrations”}
\author{}
\date{}
\linespread{1.3}
\begin{document}
	\renewcommand{\theequation}{S.\arabic{equation}}
	\setcounter{equation}{0}
	\renewcommand\thesection{S\arabic{section}}
	\renewcommand\thesubsection{S\arabic{section}.\arabic{subsection}}
	
	\maketitle


	\section*{Overview}	
		
	This document provides supplementary material for the paper titled \emph{“A Relativity-Based Framework for Statistical Testing Guided by the Independence of Ancillary Statistics: Methodology and Nonparametric Illustrations.”} This supplement presents technical proofs of results stated in the main text which are omitted there for brevity.
	
	Section~\ref{S0} 
	collects remarks that support or extend points made in the main article, including additional results.
		With respect to the results shown in Section~\ref{sc32}, Section~\ref{Sn2} provides a detailed analysis of the Anderson–Darling (AD) and Cramér–von Mises (CvM) tests applied to both the raw observations $X_1,\ldots,X_n$ and their absolute values. We examine the asymptotic relative efficiency (ARE) of these tests under a class of symmetric contiguous alternatives. The analysis, based on the framework of \citet{inglot2000vanishing}, demonstrates that applying the tests to the absolute values—i.e., folding the data—increases their asymptotic efficiency.		
	Section~\ref{S1} (“Supplementary to Section~\ref{sc33}”) presents Lemma~\ref{SL1} and Proposition~\ref{thm:ratio}, which develop the RV–based dependence analysis for Hotelling’s $T^2$ and the trace–normalized $T_1^2$ statistics defined in Section~\ref{sc33}. 
	 
	 	All references cited in this supplement are listed at the end of the document.
	
	
	\section{\ Additional Remarks}\label{S0}
	\begin{remark}\label{Sr02}
		\emph{
			Assume that a test statistic \( Q \) and a set of ancillary statistics \( V \) can be constructed from modified data \( D^M \).  
			Suppose the mapping \( D^M \to (Q, V) \) is injective with a measurable inverse.  
			If \( Q \) is independent of \( V \) under \( H_k \), \( k \in \{0,1\} \), and \( f^{Q}_1(u) = u\, f^{Q}_0(u) \),  
			then \( Q \) is MP among all test statistics that can be constructed from \( D^M \) alone.  
			For example, in several settings, it is natural to take \( D^M \) as the set of ranks of the observations; see, e.g.,~\citet{sidak1999theory,shiraishi1986optimum}.
		}
	\end{remark}
\begin{remark}\label{2.4}
	{\em
		\citet{anderson2016note} discussed criteria for evaluating and comparing the performance of test statistics. Suppose, in a given problem, we consider two test statistics, say $T_1$ and $T_2$, along with two ancillary vectors $\mathbf{V_1}$ and $\mathbf{V_2}$. It is possible that $\mathbf{V_1} = \mathbf{V_2}$. The underlying data $D$ can be reconstructed either from $(T_1, \mathbf{V_1})$ or from $(T_2, \mathbf{V_2})$. When selecting between $T_1$ and $T_2$ for use in a study, the distance correlation introduced by \citet{szekely2007measuring} between a test statistic and its corresponding ancillary vector may serve as a criterion. This approach is applicable in both parametric and nonparametric comparisons of competing test procedures.
	}
\end{remark}


\section{\ Anderson-Darling and Cram\'{e}r-von Mises Tests Based on $X$ and  $|X|$: Formulation, Symmetry, and Asymptotic Relative Efficiency.}\label{Sn2}
We begin by analyzing the Anderson–Darling (AD) tests in detail. The corresponding Cramér–von Mises (CvM) tests can be studied analogously, and we omit their derivation for brevity. Our main reference throughout is \citet{inglot2000vanishing}.

\subsection{Testing problems and statistics}

Without loss of generality, we assume that  $X_1,\dots,X_n$ are i.i.d.\ from a distribution function $F$ with $\mathrm{E}\left(X_1\right)=0$ and $\mathrm{var}\left(X_1\right)=1$. We consider
\[
H_0:\;F=\Phi, \qquad\text{vs.}\qquad H_1:\;F\neq \Phi,
\]
where $\Phi$ is the standard normal cdf. Define $U_i=\Phi(X_i)\in(0,1)$ so that under $H_0$, $U_i\sim\mathrm{Unif}(0,1)$.

\paragraph{Anderson - Darling on $X$ (test $T$).}
Let $F_n^{U}$ denote  the empirical distribution function based on the sample $U_1,\ldots,U_n$. The AD statistic is
\[
T \;=\; n\int_0^1 \frac{\{F_n^{U}(t)-t\}^2}{t(1-t)}\,dt.
\]

\paragraph{Anderson--Darling on $|X|$ (test $T_1$).}
Let $Z_i=|X_i|$ and define $Y_i = 2\Phi(Z_i)-1\in(0,1)$; under $H_0$, $Y_i\sim\mathrm{Unif}(0,1)$. 
Let  $F_n^{Y}$ denote  the empirical distribution function based on the sample $Y_1,\ldots,Y_n$.
Define
\[
T_1 \;=\; n\int_0^1 \frac{\{F_n^{Y}(v)-v\}^2}{v(1-v)}\,dv.
\]


\subsection{Local alternatives on the uniform scale}
\label{ssec:alts}
Following \citet{inglot2000vanishing}, we evaluate the efficiency of the tests under the standard class of contiguous alternatives.
\paragraph{Test $T$:} On the $U$-scale we adopt the standard contiguous alternative defined
by densities (with respect to Lebesgue measure on $[0,1]$),
\[
p_\theta(u)\;=\;1+\theta\,a(u), \qquad u\in(0,1),\quad \theta\to 0,
\]
with bounded function $a(u)$, and normalization $\int_0^1 a(u)\,du=0$ and $\int_0^1 a(u)^2\,du=1$.

%\begin{assumption}[Symmetry]
%	\label{ass:symm}
%	All alternatives are symmetric about $0$ on the $X$-scale, i.e.
%	the perturbation of the density of $X$ is even. On the $U$-scale this is equivalent to
%	\[
%	a(u)=a(1-u),\qquad 0<u<1.
%	\]
%\end{assumption}
%
%\begin{proof}[Proof that $a(u)=a(1-u)$]
%	Let $U=\Phi(X)$, so $1-U=\Phi(-X)$. Symmetry on the $X$-scale means the law of $X$ equals %that of $-X$, hence the law of $U$ equals that of $1-U$. Therefore $p_\theta(u)=p_\theta(1-u)$ %for all $u\in(0,1)$ and small $\theta$. With $p_\theta(u)=1+\theta a(u)$ we obtain %$a(u)=a(1-u)$ by cancellation. \qedhere
%\end{proof}

\begin{assumption}[Symmetry]
	\label{ass:symm}
	The alternatives are symmetric about zero on the \( X \)-scale; that is, the perturbation function \( a_X(x) \) satisfies \( a_X(x) = a_X(-x) \). This implies that on the \( U \)-scale, the score function \( a(u) \) satisfies
	\[
	a(u) = a(1 - u), \qquad u \in (0,1).
	\]
\end{assumption}

\begin{proof}[Justification of symmetry in \( a(u) \)]
	Let \( U = \Phi(X) \), so that \( 1 - U = \Phi(-X) \). Since \( X \overset{d}{=} -X \) under the symmetry assumption, it follows that \( U \overset{d}{=} 1 - U \). Therefore, the density satisfies \( p_\theta(u) = p_\theta(1 - u) \). Using the local alternative expansion \( p_\theta(u) = 1 + \theta a(u) \), we conclude \( a(u) = a(1 - u) \). \qedhere
\end{proof}


\paragraph{Test $T_1$:}

Under the alternative above, the density of \(X\) is \(f(x) = \phi(x) [1 + \theta a(\Phi(x))]\), where $\phi(x)=\exp(-x^2/2)/(2\pi)^{0.5}$. Since the alternative is symmetric, the density of \(|X|\) at \(y\) (for \(y \geq 0\)) is:
\[
g(y) = f(y) + f(-y) = \phi(y) [1 + \theta a(\Phi(y))] + \phi(-y) [1 + \theta a(\Phi(-y))].
\]
Using \(\phi(-y) = \phi(y)\) and \(\Phi(-y) = 1 - \Phi(y)\), and using  symmetry (\(a(1-\Phi(y)) = a(\Phi(y))\)), this simplifies to:
\[
g(y) = 2\phi(y) [1 + \theta a(\Phi(y))].
\]
Under the null, the density of \(|X|\) is \(g_0(y) = 2\phi(y)\).

Now, to find the density of \(Y\), we note that \(Y = 2\Phi(|X|) - 1\). The cdf of \(Y\) is:
\[
{\Pr}(Y \leq v) = {\Pr}\left(2\Phi(|X|) - 1 \leq v\right) = {\Pr}\left(|X| \leq \Phi^{-1}\left(\frac{v+1}{2}\right)\right).
\]
Let \(z = \Phi^{-1}\left(\frac{v+1}{2}\right)\). The density of \(Y\) is the derivative of the cdf:
\[
p^Y_\theta = g(z)  \left| \frac{dz}{dv} \right|.
\]
From \(z = \Phi^{-1}\left(\frac{v+1}{2}\right)\), we have \(\Phi(z) = \frac{v+1}{2}\). Differentiating both sides with respect to \(v\):
\[
\phi(z) \frac{dz}{dv} = \frac{1}{2} \implies \frac{dz}{dv} = \frac{1}{2\phi(z)}.
\]
Substituting \(g(z) = 2\phi(z) [1 + \theta a(\Phi(z))]\):
\[
p^Y_\theta(v) = 2\phi(z) [1 + \theta a(\Phi(z))] \cdot \frac{1}{2\phi(z)} = 1 + \theta a(\Phi(z)).
\]
Since \(\Phi(z) = \frac{v+1}{2}\), we get:
\[
p^Y_\theta(v) = 1 + \theta a\left( \frac{v+1}{2} \right).
\]
Thus, the alternative density for \(Y\) is 
\[p^Y_{\theta}(v) = 1 + \theta b(v),\] 
where \(b(v) = a\left( \frac{v+1}{2} \right)\).

\subsection{The MP test of $H_0$: $U\sim$ Unif$(0,1)$ (or equivalently, $Y\sim$ Unif$(0,1)$)
	against the simple alternative $U\sim p_{\theta}$ (or equivalently,  $Y\sim p^Y_{\theta}$)}
According to ~\citet[p. 224]{inglot2000vanishing}, the corresponding locally most powerful test is given by:
\[
T^{\ast}
= \bigl(\sqrt{n}\,\sigma_{0n}\bigr)^{-1}
\sum_{i=1}^n \bigl(\log p_{\theta_n}(X_i) - e_{0n}\bigr),
\qquad
e_{0n} = \operatorname{E}_{\theta_0}\!\left[\log p_{\theta_n}(X)\right],
\quad
\sigma_{0n}^2 = \operatorname{var}_{\theta_0}\!\left[\log p_{\theta_n}(X)\right]
\]
with $0<\theta_n\to 0$ as $n\to \infty$.

\subsection{Drift functionals}
\label{ssec:functionals}
As shown in \citet{inglot2000vanishing}, the Pitman slopes of the test statistics can be expressed via weighted integrals of the following score function.

\paragraph{Test $T$:} Define the integrated score on $(0,1)$
\[
A(t) \;=\; \int_0^t a(x)\,dx, \qquad 
A_w(t) \;=\; \frac{A(t)}{\sqrt{t(1-t)}}.
\]
\paragraph{Test $T_1$:} For the test based on absolute values, define
\[
B(v) \;=\; \int_0^v b(s)\,ds, \qquad 
B_w(v) \;=\; \frac{B(v)}{\sqrt{v(1-v)}}.
\]


Under symmetry we have the \emph{mapping identity}
\begin{eqnarray}\nonumber
	B(u)&=&\int_{0}^{u} a\left(\frac{x+1}{2}\right)dx=2\int_{0.5}^{(u+1)/2} a\left(x\right)dx	
	\\
	&=&2A\left(\frac{u+1}{2}\right),
	\label{Eq1}
\end{eqnarray}
since 
the condition 
\(\int_{0}^1a(x)dx=0\) and Assumption \ref{ass:symm} lead to 
\[\int_{0}^{0.5}a(x)dx+\int_{0.5}^{1}a(x)dx=0, \qquad 
\int_{0}^{0.5}a(x)dx+\int_{0.5}^{1}a(1-x)dx=0
\]
meaning that \(\int_{0.5}^1a(x)dx=0\).

\subsection{A symmetry reduction for $\|A_w\|^2$ and a comparison inequality for $\|B_w\|^2$}
Let 
\begin{eqnarray*}\nonumber
	\|A_w\|^2 \;=\; \int_0^1 \frac{A(t)^2}{t(1-t)}\,dt \qquad\mbox{and}\qquad
	\|B_w\|^2 \;=\; \int_0^1 \frac{B(t)^2}{t(1-t)}\,dt.
\end{eqnarray*}


We show that  $A(1-t)=-A(t)$.
\begin{eqnarray*}\nonumber
	A(1-u)&=&\int_{0}^{1-u} a\left(x\right)dx=\int_{0}^{1-u} a\left(1-x\right)dx=
	-\int_{1}^{u} a\left(x\right)dx=\int_{u}^{1} a\left(x\right)dx\\
	&=&-A(u),
	\label{Eq2}
\end{eqnarray*}
since 
the condition 
\(\int_{0}^1a(x)dx=0\) leads to 
\[\int_{0}^{u}a(x)dx+\int_{u}^{1}a(x)dx=0, \qquad 
\int_{0}^{u}a(x)dx+\int_{u}^{1}a(x)dx=0
\]
meaning that \(\int_{u}^1a(x)dx=-\int_{0}^ua(x)dx=-A(u)\).

Hence,
\begin{eqnarray}\nonumber
	\|A_w\|^2 \;=\; \int_0^1 \frac{A(t)^2}{t(1-t)}\,dt =
	\int_0^{0.5}\frac{A(t)^2}{t(1-t)}\,dt +\int_{0.5}^{1}\frac{A(t)^2}{t(1-t)}\,dt 
	\\
	=\int_0^{0.5}\frac{\left(A(1-t)\right)^2}{t(1-t)}\,dt +\int_{0.5}^{1}\frac{A(t)^2}{t(1-t)}\,dt 
	\;=\; 2\int_{0.5}^{1} \frac{A(t)^2}{t(1-t)}\,dt.
	\label{Eq3}
\end{eqnarray}

Thus, we conclude that
\begin{proposition}[Norm comparison]
	\label{thm:norm}
	Under Assumption~\ref{ass:symm} and for any nondegenerate $a$ (equivalently $A\not\equiv 0$),
	\[
	\|B_w\|^2 \;>\; 2\,\|A_w\|^2.
	\]
\end{proposition}
\begin{proof}
	We have that Equation~(\ref{Eq1}) gives	
	\[
	\|B_w\|^2 \;=\; \int_0^1 B_w(v)^2\,dv 
	\;=\; 4\int_{0}^{1}  \frac{A\left((t+1)/2\right)^2}{t(t-1)}\,dt\,=\,
	4\int_{0.5}^{1}  \frac{A\left(z\right)^2}{(1-z)(2z-1)}\,dz.
	\]
	For $z\in[0.5,1)$, the denominator $(1-z)(2z-1)=(1-z)(z+z-1)< (1-z)z$. Result~(\ref{Eq3})
	completes the proof.	
\end{proof}



\subsection{Asymptotic Relative Efficiency}
According to Theorem~5.2(d) of \citet{inglot2000vanishing}, we have
\[
\text{ARE}(T,T^{\ast}) = 2\|A_w\|^2 < \|B_w\|^2 = \text{ARE}(T_1,T^{\ast})/2,
\]
where $\text{ARE}$ denotes the asymptotic relative efficiency, defined as the approximate ratio of the sample sizes required by a given test and the most powerful (MP) test to attain the same power under a sequence of local alternatives $p_{\theta_n}$, with $\theta_n \to 0$ and $n \to \infty$, such that $\theta_n^\gamma \Phi^{-1}(1-\alpha_n)\to 0$ for some $\gamma \in (0,1)$ and $\alpha_n \to 0$; see \citet[p.~231]{inglot2000vanishing}.

Therefore, under symmetric alternatives around zero, the test \( T_1 \) based on the absolute values \( |X_1|, \ldots, |X_n| \) is asymptotically more efficient than the classical Anderson--Darling test \( T \), which is applied directly to the raw observations.


	\section{Supplementary to Section~\ref{sc33}}\label{S1}
\begin{lemma}\label{SL1}
	Let $\mathbf{X}_{1},\ldots,\mathbf{X}_{n}$ be i.i.d.\ $p$-dimensional random vectors with mean $\mu=\mathrm{E}(\mathbf{X}_{i})$ and covariance $\Sigma=\mathrm{var}(\mathbf{X}_{i})$. 
	Write $\bar{\mathbf{X}}=n^{-1}\sum_{i=1}^{n}\mathbf{X}_{i}$ and 
	$
	\mathbf{V}=\big(\mathbf{X}_{1}-\bar{\mathbf{X}},\ldots,\mathbf{X}_{n}-\bar{\mathbf{X}}\big)^{\!\top}\in\mathbb{R}^{n\times p}.
	$
	Let $T^{2}$ be Hotelling’s statistic (or, more generally, any statistic that is symmetric under permutations of the sample indices). Then, for each $i\in\{1,\ldots,n\}$,
	\[
	\mathrm{E}\!\left[T^{2}\,(\mathbf{X}_{i}-\bar{\mathbf{X}})\right]=\mathbf{0}\in\mathbb{R}^{p}.
	\]
	Consequently, for any $b\in\mathbb{R}^{n}$,
	\[
	\mathrm{E}\!\left[T^{2}\, b^{\top}\mathbf{V}\right]=0
	\quad\text{and}\quad
	\mathrm{Cov}\!\left(T^{2},\, b^{\top}\mathbf{V}\right)=0.
	\]
	The same conclusions hold with $T^{2}$ replaced by $T_{1}^{2}$.
\end{lemma}

\begin{proof}
	Since $T^{2}$ is a symmetric statistic of the sample $\{\mathbf{X}_{1},\ldots,\mathbf{X}_{n}\}$ and the observations are i.i.d., the vector of mixed moments $\mathrm{E}(T^{2}\mathbf{X}_{j})$ does not depend on $j$. Hence,
	\[
	\mathrm{E}\!\left[T^{2}\,\bar{\mathbf{X}}\right]
	=\mathrm{E}\!\left[T^{2}\,\frac{1}{n}\sum_{j=1}^{n}\mathbf{X}_{j}\right]
	=\frac{1}{n}\sum_{j=1}^{n}\mathrm{E}\!\left[T^{2}\mathbf{X}_{j}\right]
	=\mathrm{E}\!\left[T^{2}\mathbf{X}_{1}\right].
	\]
	Therefore, for each $i$,
	\[
	\mathrm{E}\!\left[T^{2}\,(\mathbf{X}_{i}-\bar{\mathbf{X}})\right]
	=\mathrm{E}\!\left[T^{2}\mathbf{X}_{i}\right]-\mathrm{E}\!\left[T^{2}\bar{\mathbf{X}}\right]
	=\mathrm{E}\!\left[T^{2}\mathbf{X}_{1}\right]-\mathrm{E}\!\left[T^{2}\mathbf{X}_{1}\right]
	=\mathbf{0}.
	\]
	For any $b\in\mathbb{R}^{n}$,
	\[
	\mathrm{E}\!\left[T^{2}\, b^{\top}\mathbf{V}\right]
	=\sum_{i=1}^{n} b_{i}\,\mathrm{E}\!\left[T^{2}(\mathbf{X}_{i}-\bar{\mathbf{X}})\right]
	=0.
	\]
	Since $\mathrm{E}(b^{\top}\mathbf{V})=0$ by construction, it follows that
	\[
	\mathrm{Cov}\!\left(T^{2},\, b^{\top}\mathbf{V}\right)
	=\mathrm{E}\!\left[T^{2}\, b^{\top}\mathbf{V}\right]-\mathrm{E}(T^{2})\,\mathrm{E}\!\left(b^{\top}\mathbf{V}\right)
	=0- \mathrm{E}(T^{2})\cdot 0
	=0.
	\]
	The argument for $T_{1}^{2}$ is identical, since it is also symmetric in the sample indices.
\end{proof}

\begin{proposition} [The asymptotic approximation of $R$]
	\label{thm:ratio}
	Let $\mathbf{X}_{1},\ldots,\mathbf{X}_{n}$ be i.i.d.\ $p$-dimensional random vectors with 
	$\mathrm{E}(\mathbf{X}_{i})=\mu$, $\mathrm{Cov}(\mathbf{X}_{i})=\Sigma$, and finite fourth moments. 
	Assume the elliptical fourth-moment identity
	\begin{equation}
		\mathrm{E}\!\left\{ (Y^\top B Y)(Y^\top C Y) \right\}
		= 2\,\mathrm{tr}(B\Sigma C\Sigma) + (\theta-1)\,\mathrm{tr}(B\Sigma)\,\mathrm{tr}(C\Sigma),
		\qquad Y=\mathbf{X}_1-\mu,
		\label{eq:elliptical-identity}
	\end{equation}
	for symmetric $p\times p$ matrices $B,C$, where $\theta\ge 1$ is the elliptical kurtosis factor. 
	Let
	\[
	T^{2}=n\,\bar{\mathbf{X}}^{\top} S^{-1}\bar{\mathbf{X}}, 
	\qquad 
	T_{1}^{2}=\frac{n p\,\bar{\mathbf{X}}^{\top}\bar{\mathbf{X}}}{\mathrm{tr}(S)},
	\qquad
	\psi=\frac{p\,\mathrm{tr}(\Sigma^2)}{\{\mathrm{tr}(\Sigma)\}^{2}}\ge 1,
	\]
	and denote $W=VV^\top$, where $V=[\,\mathbf{X}_{1}-\bar{\mathbf{X}},\ldots,\mathbf{X}_{n}-\bar{\mathbf{X}}\,]^{\top}$. 
	Then, as $n\to\infty$,
	\begin{equation}
		\frac{\mathrm{RV}(T_1^2,\mathrm{vec}(W))}{\mathrm{RV}(T^2,\mathrm{vec}(W))}
		=\frac{\big(2\psi+(\theta-1)p\big)^2}{\psi\,\big(2+(\theta-1)p\big)^2} + r_{n,p},
		\label{eq:ratio-main}
	\end{equation}
	where the remainder $r_{n,p}$ satisfies
	\[
	r_{n,p}=
	\begin{cases}
		O(n^{-1}), & \text{if $p$ is fixed}, \\[0.5ex]
		O(n^{-1/2}), & \text{if $p,n\to\infty$ with $p/n\to c\in(0,1)$}.
	\end{cases}
	\]
	The leading term in \eqref{eq:ratio-main} depends only on $\psi$ and $\theta$ and is free of $\mu$.
\end{proposition}

\begin{proof} We outline the main steps leading to the expansion of $R$.
	\textit{Step 1: Linearizations of $T^2$ and $T_1^2$.}  
	Let $S=\Sigma+\Delta$ with $\Delta=O_p(n^{-1/2})$ when $p$ is fixed. Expanding yields
	\begin{align}
		T^2 &= n\,\bar{\mathbf{X}}^{\top}\Sigma^{-1}\bar{\mathbf{X}}
		- n\,\bar{\mathbf{X}}^{\top}\Sigma^{-1}\Delta\Sigma^{-1}\bar{\mathbf{X}} + R_{T}, \label{eq:T2-exp}\\
		T_1^2 &= \frac{n p}{\mathrm{tr}(\Sigma)}\,\bar{\mathbf{X}}^{\top}\bar{\mathbf{X}}
		- \frac{n p}{\{\mathrm{tr}(\Sigma)\}^{2}}\,\bar{\mathbf{X}}^{\top}\bar{\mathbf{X}}\,\mathrm{tr}(\Delta) + R_{T1}, \label{eq:T1-exp}
	\end{align}
	with $R_T,R_{T1}=O_p(n^{-1})$. Thus $T^2$ depends on $\Delta$ only through $L_2(\Delta)=\mathrm{tr}(\Sigma^{-1}\Delta)$, while $T_1^2$ depends on $\Delta$ only through $L_1(\Delta)=\mathrm{tr}(\Delta)$.
	
	\textit{Step 2: Covariances with $\mathrm{vec}(W)$.}  
	Since $S=(n-1)^{-1}V^\top V$ and $W=VV^\top$, covariances between $L_1(\Delta)$ or $L_2(\Delta)$ and $\mathrm{vec}(W)$ can be expressed using fourth moments of $Y=\mathbf{X}_i-\mu$. 
	By \eqref{eq:elliptical-identity},
	\[
	\mathrm{E}\{(Y^\top BY)(Y^\top CY)\}
	=2\,\mathrm{tr}(B\Sigma C\Sigma)+(\theta-1)\,\mathrm{tr}(B\Sigma)\,\mathrm{tr}(C\Sigma).
	\]
	Applying this with $B=C$ equal to $I_p$ and $\Sigma^{-1}$, respectively, gives
	\[
	\Gamma_{I_p} = 2\,\mathrm{tr}(\Sigma^2) + (\theta-1)\{\mathrm{tr}(\Sigma)\}^2, 
	\qquad
	\Gamma_{\Sigma^{-1}} = 2p+(\theta-1)p^2.
	\]
	Hence, up to order $O(n^{-1})$, 
	\[
	\Cov(T_1^2,\mathrm{vec}(W)) \propto \Gamma_{I_p}, 
	\qquad 
	\Cov(T^2,\mathrm{vec}(W)) \propto \Gamma_{\Sigma^{-1}}.
	\]
	
	\textit{Step 3: Variances.}  
	For fixed $p$, standard quadratic-form CLTs imply
	\[
	\Var\!\left(n\,\bar{\mathbf{X}}^{\top}\Sigma^{-1}\bar{\mathbf{X}}\right) = 2p+O(n^{-1}),
	\qquad
	\Var\!\left(\tfrac{n p}{\mathrm{tr}(\Sigma)}\,\bar{\mathbf{X}}^{\top}\bar{\mathbf{X}}\right) 
	= 2p\,\psi+O(n^{-1}).
	\]
	For proportional growth $p/n\to c\in(0,1)$, both variances inflate by a factor of $(1-c)^{-3}$, and fluctuations scale as $O(n^{-1/2})$.
	
	\textit{Step 4: RV ratios.}  
	By the definition of RV,
	\[
	\mathrm{RV}(U,\mathbf{Y})=\frac{\|\Cov(U,\mathbf{Y})\|^2}{\Var(U)\,\Var(\mathbf{Y})}.
	\]
	Since the common factor $\Var(\mathrm{vec}(W))$ cancels in the ratio, combining Steps 2 and 3 gives
	\[
	\frac{\mathrm{RV}(T_1^2,\mathrm{vec}(W))}{\mathrm{RV}(T^2,\mathrm{vec}(W))}
	=\frac{\{\Gamma_{I_p}\}^2/\Var(T_1^2)}{\{\Gamma_{\Sigma^{-1}}\}^2/\Var(T^2)}+r_{n,p}.
	\]
	Substituting $\Gamma_{I_p},\Gamma_{\Sigma^{-1}}$ and the variances yields
	\[
	\frac{\mathrm{RV}(T_1^2,\mathrm{vec}(W))}{\mathrm{RV}(T^2,\mathrm{vec}(W))}
	=\frac{\big(2\psi+(\theta-1)p\big)^2}{\psi\,\big(2+(\theta-1)p\big)^2}+r_{n,p},
	\]
	with $r_{n,p}=O(n^{-1})$ for fixed $p$, and $r_{n,p}=O(n^{-1/2})$ for proportional growth $p/n\to c\in(0,1)$.
	
	\textit{Step 5: Independence of $\mu$.}  
	The leading term in \eqref{eq:ratio-main} does not involve $\mu$, since $\mu$ affects only the mean of $\bar{\mathbf{X}}$ while the quadratic-form variances and fourth-moment identities depend solely on centered quantities.
	
\end{proof}




%%%%%%%%%%%%%%%%%%%%%%%%%%%%%%%%%%%%%%%%%%%%%%%%%%%%%%%%%%%%%%%%%%		

	

	
\bibliographystyle{Chicago}
\bibliography{testsApproachREF.bib}
\end{document}

